% ============================================================
% 第5章 标签构造、样本划分与评估方案
% 【负责人：验证负责人（已按统一配置预填“数据侧承诺”，折分与指标
%   细节以验证负责人冻结版为准）】
% ============================================================
\section{标签构造、样本划分与评估方案}

\subsection{标签定义}

\begin{itemize}
  \item \textbf{正类}：标签窗口（2026-06-21 起 40 天）内出现至少一条有效的
        11803（事故）或 11804（未遂事故）记录。主口径实测：正类车
        \textbf{126} 台 / 标签记录 \textbf{222} 条（清洗前后守恒）。
  \item \textbf{负类}：窗口内无 11803/11804 记录即记 $y=0$。团队决策
        （2026-09-18）：事件日志视为完整，不再区分“确认负类”与“代理负类”；
        若官方日后说明日志不完整，回滚此假设并重算全部基线。
  \item \textbf{清洗与标签的关系}：坐标无效只置字段缺失并标记，不得删除行或车
        ——主口径标签窗内 10 条 $(0,0)$ 记录（含事故样本车辆A）全部保留
        （见 3.4 节核心纪律）；事故 11803 与未遂 11804 \textbf{不套用}普通报警的
        30 秒段合并规则。
\end{itemize}

\subsection{样本划分：车辆级分层 5 折}

\begin{itemize}
  \item 车辆级划分（同车全部窗口在同一折），按标签分层，保证每折正类
        25--26 台；折号由哈希确定，完全可复现，登记于 \texttt{folds.csv}
        （含 \texttt{gpsno}$\to$\texttt{fold} 映射与 \texttt{split\_version}）。
  \item 填补、标准化、平滑、分箱、PCA、相关聚类、特征选择、校准、阈值与融合
        \textbf{均只在训练折拟合}。
  \item 每次实验保存每车折外预测、折号、种子与配置版本，保证任何数字可回溯。
\end{itemize}

\subsection{评价指标体系}

以 AUC 为唯一主指标（与官方一致），配套指标用于运营性判断：

\begin{table}[H]
\centering
\caption{评价指标体系}
\label{tab:metrics}
\small
\begin{tabularx}{\textwidth}{lL}
\toprule
指标 & 用途 \\
\midrule
\texttt{roc\_auc\_overall\_oof} & 主指标：500 台车折外预测的整体 AUC \\
\texttt{roc\_auc\_per\_fold} & 折间稳定性（5 折逐一报告） \\
average\_precision & 正类稀缺下的排序参考（非梯形 PR AUC） \\
lift@10\% / lift@20\% & 预警名单的运营价值 \\
recall@K / precision@K & 管理容量 K 下能覆盖多少高风险车 \\
brier\_score / log\_loss & 概率校准质量 \\
\bottomrule
\end{tabularx}
\end{table}

\subsection{实验准入规则（全组统一纪律）}

折外 AUC 提升 \textbf{0.005} 为提示线，四条同时满足才进入主线：
\begin{enumerate}[label=(\arabic*)]
  \item 至少 4/5 折不下降；
  \item 车辆级配对 bootstrap 区间下界 $\ge -0.005$（同一车辆新旧预测成对重采样）；
  \item 折间方向一致；
  \item 不伤害缺失子群（如 IMU 缺口车辆）。
\end{enumerate}

\begin{TODO}
种子集合（12/34/56/78/90）下的多种子均值与方差、阈值选择规则（官方未给则按
管理容量 K 或召回目标在训练内层选择）请验证负责人补充冻结说明与最终版本号。
\end{TODO}
